\section{The full intuitionistic theory and both composites}
\label{sec:bridge}

Exact equivalence of clause sequents is stronger than equality of closed
equational theorems. Here we spell out how it transfers the complete
iPRA interpretation, including its logical operations and both round
trips. This section summarizes inherited results~\cite{repo,lean,arithmetic};
their detailed written arithmetic proofs are included in \texttt{reference}.

\subsection{Ordinary realization}
Let $T_I$ be intuitionistic first-order logic with actual equality,
genuine falsity $\bot$, and the staged scalar-iterator language. Its
nonlogical axioms and schemas are the two iterator computations,
\[
 \forall x\,(x=0\vee\exists y\,x=Sy),
 \qquad (0=S0)\to\bot,
\]
and induction for every literal equation $E(y,\vec p)$:
\[
 \bigl(E(0,\vec p)\wedge
       \forall y(E(y,\vec p)\to E(Sy,\vec p))\bigr)
       \to\forall y E(y,\vec p).
\]
Equivalently this is the usual contextual induction rule, with the
induction variable fresh for retained assumptions. A staged iterator
symbol $i_h(n,a,\vec p)$ is introduced for each fixed earlier-stage
description $h(u,\vec p)$; no object-language variable ranges over
arbitrary function descriptions. The numerical domain is the whole
domain and interpreted equality is equality itself.

Use has its ordinary conditional-universal reading. Introduce is sound
by numerical coverage in case mode and equation induction in progress
mode. Opening is sound by Lemma~\ref{lem:opening-sound}. The explicit
falsity bridge must not be omitted from $T_I$: positive arithmetic
equations alone admit a one-element structure. In the native calculus
the designated contradiction is the equation $0=S0$.

\subsection{All logical operations as macros}
Write $C(v,a,b)=N(I(v,a;u.b))$ with $u$ unused, and let
$q_y A(y)=D_{w,y}A(y)$ for unused $w$. A gate on formulas is defined by
external recursion on their finite syntax:
\[
 \begin{aligned}
 G_v(s=t)&=(C(v,0,s)=C(v,0,t)),\\
 G_v(D_{x,y}B)&=D_{x,y}G_vB.
 \end{aligned}
\]
All bound variables are renamed away from $v$. Its endpoints satisfy
$\vdash G_0A$ and $G_{Sz}A\dashv\vdash A$; also
$A\dashv\vdash q_yG_yA$ when $y$ is fresh. These are derivability
statements, including for nested binders.

For equations define the selector
\[
 M_v(s=t,u=w)=(C(v,s,u)=C(v,t,w)).
\]
For compound formulas align fresh paired binders and define
\[
 \begin{aligned}
 M_v(D_{x,y}A,D_{x,y}B)&=D_{x,y}M_v(A,B),\\
 M_v(E,D_{x,y}B)&=D_{x,y}M_v(G_yE,B),\\
 M_v(D_{x,y}A,E)&=D_{x,y}M_v(A,G_yE),
 \end{aligned}
\]
where $E$ is an equation. Each recursive call lowers the maximum binder
depth. Gate padding and binder congruence give
$M_0(A,B)\dashv\vdash A$ and $M_{Sz}(A,B)\dashv\vdash B$.
With $P(y)=N(I(y,0;u.C(u,1,0)))$, the following complete inventory is
available:
\[
\begin{aligned}
 \bot_*&=(0=S0),\\
 A\Rightarrow_*B&=q_yM_y(A,B),\\
 A\wedge_*B&=q_yG_yM_{P(y)}(A,B),\\
 \exists_*x\,A(x)&=D_{x,y}G_yA(x),\\
 A\vee_*B&=\exists_*v\,M_v(A,B),\\
 \forall_*x\,A(x)&=A(0)\wedge_*q_yG_yA(y).
\end{aligned}
\tag{Logical macros}
\]
Freshness is understood in every line. The last line intentionally
supplies the zero instance separately: the gated closure supplies all
successor instances.

Implication introduction and elimination use selector endpoints and the
closure rules. Existential introduction and elimination use gate
endpoints. Conjunction is projected at $1$ and $2$, using $P(1)=1$,
$P(2)=0$; introduction uses numerical cases on its selector. Disjunction
uses control witnesses $0$ and $1$, and its elimination uses numerical
cases. For universals, numerical coverage combines the zero conjunct
with the successor instances in the closure. Under $0=S0$, equality
transport identifies $C(0,s,t)$ and $C(S0,s,t)$, proving every numerical
equation; external induction on formula depth and case introduction
then prove every formula. Thus all usual intuitionistic first-order
inference rules, including explosion, are recovered.

These macros do not add unrestricted induction for quantified formulas.
Primitive progress remains restricted to equations. Equation induction itself is obtained
by introducing a closure at an unused zero witness, opening it, and
supplying its base. The explicit repaired case rule handles the case
introduction instances needed above.

\subsection{The numerical arithmetic bridge}
The inherited scalar proof has a necessary dependency order. First it
constructs a scalar predecessor and proves $\operatorname{pred}(0)=0$
and $\operatorname{pred}(Sn)=n$ using equation induction and coverage.
It then constructs truncated subtraction and a characteristic term
for actual equality on the entire domain. Compiling a fixed
quantifier-free formula into its characteristic equation gives
quantifier-free induction. Only after this stage does it construct
pairing and compile arbitrary primitive recursion. These steps are
proved internally, rather than checked only on standard numerals.

For concreteness, once predecessor, subtraction, and quantifier-free
induction are available, a bounded square-root construction proves
pairing and projections $\pi,L,H$ satisfying
\[
 L(\pi(a,b))=a,\qquad H(\pi(a,b))=b.
\]
One choice is $\pi(a,b)=(a+b)^2+a$; bounded search for the integer
square root supplies its projections. Surjectivity of $\pi$ is
unnecessary: it stores a recursion state and does not replace the
interpretation domain by a coding range.

Given a fixed primitive recursion
\[
 f(0,\vec p)=g(\vec p),\qquad
 f(Sn,\vec p)=h(n,f(n,\vec p),\vec p),
\]
and already translated $g^*,h^*$, set
\[
 \begin{aligned}
 W(u,\vec p)&=\pi(SL(u),h^*(L(u),H(u),\vec p)),\\
 V(n,\vec p)&=I(n,\pi(0,g^*(\vec p));u.W(u,\vec p)),\\
 f^*(n,\vec p)&=H(V(n,\vec p)).
 \end{aligned}
\]
Equation induction proves $L(V(n,\vec p))=n$ and the defining equations
for $f^*$. Initial functions and composition translate directly.
In the reverse direction, every fixed scalar iterator is an ordinary
primitive recursion whose step ignores the recursion counter.

\subsection{Why this is bi-interpretability}
For the numerical composites, internal equation induction proves
uniqueness of each fixed iterator or primitive recursion satisfying
its defining equations. External induction over a finite function
description therefore identifies its double translation with the
original operation at every internal argument. The identity graph
$x=y$ is consequently a definable total bijection preserving every
operation in both composite interpretations. It is not a retract,
quotient, or restriction to standard numerals.

For the logical composites, the essential native identity is
\[
 D_{x,y}B(x,y)\dashv\vdash
 \exists_*x\bigl(B(x,0)\Rightarrow_*\forall_*y B(x,y)\bigr).
\]
Forward, ordinary opening supplies the conditional-universal clause;
introduce the three macros and close the opening. Backward, eliminate
the existential macro, introduce $D$ at its witness using the repaired
case rule, and obtain the required successor instance from implication
and universal elimination. The endpoint proofs above give the round
trips for all ordinary connectives and quantifiers. External induction
on formulas finishes both composites. Optional clauses correspond to
their unconditional or conditional universal readings.

These are exactly the inherited bridges recorded by
\texttt{biInterp} and the four faithfulness and formula round-trip theorems in
\texttt{NativeBiInterpretation.lean}~\cite{lean}. The new theorem
preserves every end sequent of the canonical clause presentation,
so the interpretations and the identity numerical graphs transfer
unchanged. In particular the new calculus derives $0=S0$ if and
only if iPRA derives $\bot$. This establishes equiconsistency; it
does not assert that either theory proves its own consistency.
